\chapter{绪论}
\label{ch:introduction}

\section{课题背景及研究目的与意义}
\subsection{研究背景}

近年来，随着低空经济与智能感知技术的发展，无人机（Unmanned Aerial Vehicle，UAV）在环境监测、设施巡检和应急救援等领域得到了广泛应用\cite{shakhatreh2019civil}。2024 年，工业和信息化部等四部门联合印发《通用航空装备创新应用实施方案（2024--2030 年）》，明确提出拓展无人机行业应用并推进应急救援示范\cite{miit2024aviation}。在灾后侦察等复杂作业中，无人机能够进入人员难以抵达的区域开展观测，但受地面通信基站受损和机载电池容量有限等因素制约，现场信息往往难以持续、及时地回传至后方\cite{erdelj2017disaster}。

面对大范围、地形复杂的受灾区域，单架无人机在作业效率和视野覆盖上均难以满足应急需求，实际任务中通常部署多架无人机协同开展观测。多机协同一方面能够分区域并行排查，显著缩短搜寻耗时；另一方面，对于存在遮挡或结构复杂的局部目标，多架无人机从不同方位同步观测能够形成视角互补，降低单一视角造成的误判风险。以图~\ref{fig:intro:scenario} 所示的灾后救援通道快速评估为例，当强降雨或地震导致某段道路周边出现滑坡、建筑倒塌和积水时，救援车队正在前方待命，指挥中心需要在极短的应急调度窗口内判断该路段是否具备重型车辆通行条件。此时，受倒塌物、桥体或植被遮挡，单架无人机从一侧低空往往只能看清局部路面状况，无法同时确认另一侧路基是否掏空、上方是否存在危石威胁或周边有无绕行空间，直接依据单机画面下结论极易误判；如果仅靠一架无人机逐段往返绕飞并反复调整角度取证，又会耗费过多时间与电量，难以赶在调度时限内完成大范围评估。因此，需要多架无人机从不同视角协同观测，并将局部观测信息经空中中继发回地面中心综合研判。

\begin{figure}[H]
\centering
% Illustrative application background, not an experimental deployment.
% This scene explains the observation and communication requirements;
% algorithm relationships are given separately in research-framework.tex.
\begin{tikzpicture}[
  x=1mm,y=1mm,>=Stealth,
  font=\fontsize{9bp}{11.5bp}\selectfont,
  bgnote/.style={align=center,fill=white,inner sep=.6mm},
  bgview/.style={draw=black!35,line width=.4pt,densely dotted},
  bgdata/.style={->,draw=black!85,line width=.8pt},
  bgfeedback/.style={->,draw=black!75,line width=.65pt,
    dash pattern=on 2.5pt off 2pt},
  pics/bgdrone/.style={code={
    \draw[black!85,line width=.7pt]
      (-3.4,-1.6)--(3.4,1.6) (-3.4,1.6)--(3.4,-1.6);
    \foreach \px/\py in {-4.1/-1.9,-4.1/1.9,4.1/-1.9,4.1/1.9}
      \draw[black!80,fill=white,line width=.6pt]
        (\px,\py) ellipse[x radius=1.8mm,y radius=.75mm];
    \draw[black!85,fill=black!18,line width=.6pt]
      (0,0) ellipse[x radius=1.45mm,y radius=1mm];
    \draw[black!80,line width=.5pt] (-.8,-.7)--(-.8,-1.6)
      (.8,-.7)--(.8,-1.6);
    \draw[black!80,fill=white,line width=.45pt]
      (-.55,-1.7) rectangle (.55,-2.5);
  }}
]
\path[use as bounding box] (0,0) rectangle (148,88);
\node[anchor=west,font=\fontsize{9bp}{11.5bp}\selectfont\bfseries]
  at (1,85.5) {灾后道路观测（应用示例）};

% Road and the partially overlapping fields of view, in oblique projection.
\filldraw[fill=black!2,draw=black!35,line width=.45pt]
  (3,27)--(91,27)--(110,46)--(22,46)--cycle;
\filldraw[fill=black!12,draw=black!45,line width=.45pt]
  (8,31)--(96,31)--(104,39)--(16,39)--cycle;
\draw[white,line width=.9pt,dash pattern=on 3mm off 2mm]
  (12,35)--(100,35);
\fill[black,opacity=.025]
  (20,62.5)--(7,29)--(46,29)--(61,44)--(22,44)--cycle;
\fill[black,opacity=.025]
  (66,58.5)--(44,29)--(90,29)--(105,44)--(59,44)--cycle;
\draw[black!45,line width=.45pt,densely dotted]
  (7,29)--(46,29)--(61,44)--(22,44)--cycle
  (44,29)--(90,29)--(105,44)--(59,44)--cycle;
\foreach \p in {(7,29),(46,29),(22,44)}
  \draw[bgview] (20,62.5)--\p;
\foreach \p in {(44,29),(90,29),(105,44)}
  \draw[bgview] (66,58.5)--\p;

% A building occludes part of the road from the left-hand viewpoint.
\filldraw[fill=black!10,draw=black!65,line width=.5pt]
  (42,38)--(52,38)--(52,47)--(42,47)--cycle;
\filldraw[fill=black!23,draw=black!65,line width=.5pt]
  (52,38)--(58,44)--(58,53)--(52,47)--cycle;
\filldraw[fill=black!5,draw=black!65,line width=.5pt]
  (42,47)--(52,47)--(58,53)--(48,53)--cycle;
\node[font=\fontsize{8bp}{10bp}\selectfont] at (47,42.5) {建筑};

% Water on the road is a concrete observation target, with no claimed result.
\filldraw[fill=black!22,draw=black!55,line width=.45pt]
  (67,33.2) .. controls (65,35) and (69,37) .. (73,37)
  .. controls (77,38.8) and (85,37.3) .. (86,35.4)
  .. controls (89,33.4) and (84,32) .. (79,32.5)
  .. controls (74,31.8) and (70,32) .. cycle;
\draw[white,line width=.65pt]
  (70,34.5) .. controls (73,35.5) and (76,33.8) .. (79,34.8)
  (77,36) .. controls (79,36.7) and (82,35.5) .. (84,36);
\node[bgnote,inner sep=.3mm] at (24,28.1) {道路入口};
\node[bgnote,inner sep=.3mm] at (78,28.1) {中段积水};
\node[font=\fontsize{8.5bp}{10.5bp}\selectfont,align=center]
  at (53,22.5) {不同视角提供局部观测，部分区域重复覆盖};

% Observing UAVs and the aerial relay.
\pic at (20,65) {bgdrone};
\pic at (66,61) {bgdrone};
\pic at (92,74) {bgdrone};
\node[bgnote] at (20,56.6) {观测无人机 A};
\node[bgnote] at (76,52.8) {观测无人机 B};
\node[bgnote] at (92,82) {中继无人机};

% Data travels toward the receiver; reception feedback travels back.
\draw[bgdata] (26.5,66.4)
  .. controls (44,74) and (64,76) .. (85.2,74.6);
\node[bgnote] at (51,75.7) {观测信息回传};
\draw[bgdata] (71.7,63.6)--(86.3,71.5);
\draw[bgdata] (98.5,72.5)--(123,54.5);
\node[bgnote,inner sep=.5mm] at (112,66.5) {信息转发};
\draw[bgfeedback] (139,55)
  .. controls (150,72) and (122,84) .. (98.8,77);
\node[bgnote,inner sep=.5mm] at (130,77.5) {接收确认};

% The ground receiver combines delivered observations for the task decision.
\draw[black!80,fill=white,line width=.7pt]
  (121,44) rectangle (135,54);
\draw[black!40,line width=.4pt] (122.2,46)--(133.8,46);
\draw[black!80,line width=.65pt]
  (128,44)--(128,41.5) (124,41.5)--(132,41.5);
\draw[black!65,line width=.5pt]
  (123,51.8)--(126.5,51.8) (123,49.5)--(130,49.5)
  (131.5,51.8)--(133,51.8);
\draw[black!80,line width=.6pt] (138.5,44)--(138.5,53.3);
\fill[black!80] (138.5,53.7) circle[radius=.6mm];
\draw[black!65,line width=.5pt]
  (139.8,52.5) arc[start angle=-40,end angle=40,radius=1.8mm]
  (140.8,51.7) arc[start angle=-40,end angle=40,radius=3mm];
\node[bgnote] at (128.5,36.8) {地面接收中心};
\node[align=center,font=\fontsize{9bp}{11.5bp}\selectfont]
  at (128.5,26.5) {\textbf{任务目标}\\在时限内汇总信息\\正确判断道路状况};

% Background constraints and their consequences, separate from the methods.
\draw[black!45,line width=.45pt] (1,17.5)--(147,17.5);
\draw[black!25,line width=.4pt] (49.5,6)--(49.5,15.5)
  (98.5,6)--(98.5,15.5);
\node[align=center] at (25,11)
  {\textbf{电量有限、链路变化}\\持续回传受到限制};
\node[align=center] at (74,11)
  {\textbf{视野受限、观测重叠}\\全量上传占用带宽};
\node[align=center] at (123,11)
  {\textbf{数据丢失、反馈延迟}\\按时完成判断更困难};
\node[align=center,font=\fontsize{8bp}{10bp}\selectfont,text=black!75]
  at (74,1.5) {点线：观测范围\qquad 实线箭头：观测信息\qquad 虚线箭头：接收确认};
\end{tikzpicture}

\bicaption{通信资源受限的多无人机协同观测任务示意}{Multi-UAV collaborative observation under communication resource constraints}
\label{fig:intro:scenario}
\end{figure}

在这类灾后排查任务中，道路能否通行以及后续如何调度，最终仍要由后方地面中心汇总各方观测结果后决定。前端各架无人机仅掌握局部画面且机载算力有限，难以替代地面中心完成多视角信息的综合判断。由于灾区地面通信基站往往损毁或缺乏覆盖，距离较远的无人机无法与地面中心直接通信，只能依靠彼此组建多跳无线网络，经由邻近无人机接力转发，将现场信息回传至地面中心\cite{zeng2016uav}。但在实际作业中，这种空中多跳回传通道很难长期稳定维持。无人机在飞行过程中相对位置不断变化，当承担中继任务的节点移出通信范围时，原有转发路径便会中断。同时，随着参与观测的无人机数量增加，各路回传数据往往集中流向少数靠近地面中心的中继节点。多条链路在相近区域共享无线信道并发传输，会加剧信号间的相互干扰，并使关键中继承受过重的转发负荷。当到达的数据量超过链路转发能力时，待发数据会在中继缓存中不断积压，引起明显排队延迟甚至溢出丢包；持续高负荷转发还会加速消耗这些节点的机载电量，使其提前退出网络，进而导致依赖该中继转发的其他无人机失去回传路径。

如果前端无人机直接把高清原始图像或连续视频经上述多跳网络回传，庞大的数据量会进一步加重中继负担，迅速耗尽有限的无线带宽与机载电量，因此需要在发送前先在本地压缩待传信息。近年来，视觉语言模型（Vision-Language Model，VLM）的轻量化与端侧部署技术逐步成熟，使无人机在机载平台上直接解析观测画面成为可能。例如，MiniCPM-V 等模型借助高效视觉编码、量化及推理优化，已支持在移动设备上完成图像理解与文本生成\cite{yao2025minicpmv}；已有研究将量化后的视觉语言模型部署至无人机机载计算平台，在真实飞行中实现了基于图像与语言指令的目标定位\cite{tayar2026vlnonthefly}；面向电力巡检场景，Farahani 等在 Jetson 嵌入式平台上运行视觉语言模型，自动生成结构化报告并仅保留必要证据帧，有效减少了原始视频的连续回传\cite{farahani2026evlm}。上述研究表明，无人机可以依托机载算力在本地直接提取与当前任务相关的语义信息，并将观测画面转化为简短的类别判断、核心陈述或附带依据的详细摘要等不同详略程度的文本消息，从而避免持续回传高开销的原始图像与视频。

基于上述端侧语义提取方式，图~\ref{fig:intro:task-flow} 给出了面向任务的多无人机协同观测信息处理流程\cite{zhou2024goal}，包括局部采集、消息生成与选择、无线回传、接收融合及任务判断等环节。任务执行时，各架无人机先采集目标区域的局部画面，利用机载视觉语言模型针对当前排查问题，生成从简短类别结论到详细描述等不同长度的候选消息。由于单次任务允许传输的数据总量有限，系统不会将所有候选全部回传，而是结合通信预算从中挑选最有价值的观测视角与消息内容。选定的消息随后经无线直连或空中多跳中继发往地面中心。考虑到无线信道存在丢包，地面中心收到消息后会发回确认信号，发送端结合收到的确认情况，在规定时限内动态决定继续发送新消息、重传未确认消息、暂时等待还是停止发送。最终，地面中心汇总截止时间前实际收到的全部消息，融合多视角信息得出最终的判断结果。

在灾后道路排查中，现场险情不断变化，救援队伍也在等待前方指令，每项排查任务都有严格的截止时间。观测消息如果超过时限才送达，即便内容完整也已错过调度时机；就算按时送达，只要其中来自遮挡视角的信息存在偏差，或者缺少其他关键角度的补充，地面中心依然可能判断出错。而当已收到的信息已经足够得出准确结论时，继续发送重复内容只会额外消耗带宽和电量。因此，本文将“地面中心在规定时限内收到有效观测信息并作出正确判断”作为单次任务的完成标准，并以任务按时正确完成率来衡量系统的整体表现。

\begin{figure}[H]
\centering
% Conceptual observation process, not an integrated implementation of the
% three proposed policies. Feedback returns to transmission control.
\begin{tikzpicture}[
  x=1mm,y=1mm,>=Stealth,
  font=\fontsize{8.8bp}{11.5bp}\selectfont,
  taskbox/.style={draw=black,line width=.5pt,align=center,
    text width=22mm,minimum height=19mm,inner sep=1mm,outer sep=0pt},
  taskedge/.style={-{Stealth[length=1.6mm,width=1.1mm]},line width=.55pt},
  tasknote/.style={align=center,inner sep=0pt}
]
\path[use as bounding box] (0,0) rectangle (148,-38);

\node[taskbox] (observe) at (12,-24)
  {\textbf{局部观测}\\多节点观测\\互补或重叠};
\node[taskbox] (message) at (40,-24)
  {\textbf{生成与选择}\\确定信息来源\\与表达长度};
\node[taskbox] (transmit) at (68,-24)
  {\textbf{发送与转发}\\安排路径功率\\及发送时机};
\node[taskbox] (receive) at (96,-24)
  {\textbf{接收与融合}\\汇集已收消息\\形成任务判断};
\draw[taskedge] (observe.east) -- (message.west);
\draw[taskedge] (message.east) -- (transmit.west);
\draw[taskedge] (transmit.east) -- (receive.west);

\draw[taskedge,dashed] (receive.north) -- (96,-3.5)
  -- node[fill=white,inner sep=1mm] {接收确认} (68,-3.5)
  -- (transmit.north);
\node[tasknote,font=\fontsize{8.3bp}{10.5bp}\selectfont]
  at (82,-9.5) {可能迟到或丢失};
\draw[taskedge] (receive.east) -- (116,-24);

\draw[line width=.5pt] (116,-10) rectangle (148,-37);
\node[tasknote,font=\fontsize{9bp}{11.5bp}\selectfont\bfseries]
  at (132,-14) {任务可靠完成};
\draw[line width=.35pt] (118,-17) -- (146,-17);
\node[tasknote,font=\fontsize{8bp}{10bp}\selectfont]
  at (132,-20.5) {同时满足以下条件};
\node[tasknote] at (132,-25) {收到观测信息};
\node[tasknote] at (132,-29.5) {任务判断正确};
\node[tasknote] at (132,-34) {未超过任务时限};
\end{tikzpicture}

\bicaption{多无人机协同观测的信息流与任务完成条件}{Information flow and task completion conditions in multi-UAV collaborative observation}
\label{fig:intro:task-flow}
\end{figure}

在实际回传过程中，系统同时受到通信能量、单次任务消息总长度和任务时限三方面条件的制约。机载通信电量决定了无人机网络能够持续服务多长时间，且观测任务往往连续到达，当前任务耗电过多会直接减少后续任务的可用电量。单次任务允许发送的消息总长度有限，包含更多细节的长消息虽然信息量更大，但在相同链路速率下传输耗时更长，中途因信道波动或链路变化发生丢包的风险也更高。任务时限限制了发送数据、等待确认和丢包补发的总时间。如果为了加快传输而提高发射功率，节点能耗和对邻近链路的信号干扰都会随之上升，而丢包后的反复重传又会同时占用剩余时间与电量。为了在上述约束下提高任务按时正确完成率，多无人机协同观测回传需要解决以下三个问题。

（1）机载通信电量有限与网络拓扑高动态变化制约了信息的持续回传。无人机通过无线链路组成飞行自组网（Flying Ad Hoc Network，FANET），节点高速运动会导致邻居关系和可用路径频繁变化\cite{bekmezci2013fanet,intro4}。节点移出通信范围或关键中继电量耗尽都会造成路径中断。在功率与路径控制上，长期采用高发射功率或将流量集中在少数位置有利的中继上，会加速关键节点耗电并缩短网络服务时间；如果一味降低功率或强行均摊转发负载，又可能因速率不足或经由劣质链路导致数据积压与丢包\cite{mozaffari2019tutorial,chen2022uav}。因此，需要结合本地数据队列、实时信道状态与邻居剩余电量联合决策下一跳和发射功率，在维持较高转发吞吐量的同时延长网络可用时间。

（2）有限的消息长度预算使多视角信息的充分性与通信开销难以兼顾。多架无人机的视野之间既存在互补，也难免出现区域重叠，全部回传会迅速耗尽有限预算，过度压缩长度又容易丢失判断所需的关键细节。尽管现有协同感知研究已尝试通过筛选信息来源和空间区域来削减冗余\cite{who2com,where2comm,how2comm}，但在基于端侧视觉语言模型的语义观测中，受遮挡视角同样可能输出高置信度的错误判断，篇幅较长的详细描述也未必比简短的类别判断更具性价比；同时，地面中心每收到一条新消息，剩余候选内容的补充价值就会随之改变。因此，需要在有限消息长度预算内，综合各视角的观测可靠性、地面中心已有判断与链路传输代价，动态决定消息来源、详略程度以及停止发送的时机。

（3）严格的任务时限与不完备的确认反馈使关键消息难以按时可靠交付。选出的消息在无线信道传输中可能因信道波动而丢失，但地面中心返回的确认消息本身也可能延迟到达甚至丢失。当发送端未能按时收到确认时，无法直接分辨消息究竟是传输失败，还是已经送达接收端但确认尚在途中。此时如果立刻重传，极易对已送达消息重复发送并浪费电量与时间；如果一直等待确认，真正丢失的关键消息又可能错过补发窗口而导致任务超时。传统自动重传请求（Automatic Repeat Request，ARQ）协议主要依据确认和固定超时机制处理分组丢失\cite{lin1984arq}，未能结合协同观测的判决结果来判断补发是否必要。在实际任务中，漏掉尚未送达的互补证据可能导致中心误判，而重复补发无法改变最终类别判断的消息则没有实际作用。因此，发送端需要依据发送记录与累计确认动态推断接收状态，在剩余时间与电量约束下协调发送新消息、重传、等待与停止动作。

上述三项挑战分别对应多机协同观测回传中的网络转发、内容筛选与可靠交付环节，共同决定了地面中心能否在规定时限内作出准确判断。本文以灾后道路排查为典型场景，重点研究在通信能量、消息长度预算与任务时限共同受限的条件下，提升多无人机协同观测任务按时正确完成率的方法。

\subsection{研究目的与意义}

针对通信电量、单次任务消息长度和任务时限均受限的多无人机协同观测场景，本文的研究目的是改进多跳网络转发、多视角语义筛选以及无线可靠交付三个环节的通信策略，在控制通信耗电与传输开销的前提下，使地面中心能够在规定时限内收到足够的有效观测信息并作出准确判断，从而提高任务按时正确完成率。

从技术发展来看，现有无人机自组网研究多围绕链路吞吐量、投递率等底层指标展开，没有区分不同观测消息对最终任务判断的价值差异；而协同感知与语义通信研究又往往假设底层链路连通稳定或确认反馈及时，难以直接应对多跳中继过载、确认迟到和跨任务电量耗尽等实际问题。本文将底层路径与功率控制、前端多视角消息筛选以及传输层重传等待决策统一到同一任务完成标准下，既克服了单纯追求底层吞吐而忽视内容价值的局限，也补齐了端侧语义消息在动态无线网络中可靠回传的短板。

在实际应用方面，灾后道路排查、偏远区域设施巡检和野外应急搜救等任务大多处在无地面基站覆盖的环境中，无人机集群不仅需要长时间维持空中通信链路，还要在有限时间内把关键现场情况准确传回后方。合理调度中继功率与路径、按需筛选互补的语义描述并抑制无效重传，能够在不增加机载硬件与电池负担的前提下，延长集群的连续作业时间并提高后方的判断准确度，为无网或弱网环境下的多无人机协同观测系统建设提供直接的技术支持。

\section{国内外研究现状}
\label{intro:related}

多无人机协同观测需要将分散的信息送至接收中心，并在通信资源约束下支持任务判断。与这一过程直接相关的研究包括路由与资源分配、协同观测中的消息选择、面向任务的语义通信，以及消息丢失后的重传控制。以下按照各类方法解决的问题介绍研究进展，最后归纳本文需要进一步处理的不足。

\subsection{无人机网络路由与通信资源分配}

路由首先需要在不断变化的连接中找到可用路径。自组网按需距离矢量路由（Ad Hoc On-Demand Distance Vector，AODV）在需要通信时发现路径，减少无业务期间的维护开销\cite{aodv}；优化链路状态路由（Optimized Link State Routing，OLSR）主动交换拓扑信息，并用多点中继减少控制消息的重复转发\cite{olsr}。两者分别适合按需建立和持续维护路径，但节点移动较快时，都需要处理链路失效和控制信息过期。贪婪周边无状态路由（Greedy Perimeter Stateless Routing，GPSR）利用邻居和目的节点的位置进行局部转发\cite{gpsr}，降低了维护全局拓扑的需求。地理位置能够反映转发进度，却不能单独反映链路质量和中继电量，因而仍需结合通信状态选择路径。

无人机路由在此基础上进一步考虑运动和能量。Lakew 等归纳了飞行自组网中拓扑变化、路径维护与网络寿命之间的关系\cite{intro4}。基于 Q 学习的地理路由方法 QGeo 利用交互结果更新邻居的转发价值\cite{qgeo}；基于 Q 学习的拓扑感知路由（Q-learning-based Topology-Aware Routing，QTAR）将拓扑状态纳入下一跳选择\cite{qtar}。这类方法通过历史通信结果适应环境变化，其选择偏好由状态和回报共同决定。Li 等提出基于强化学习的节能快速路由，结合电池、信道及一跳邻居状态选择转发节点和功率，并同时考虑投递结果、时延与能耗\cite{drlfr}。Jin 等则面向城市应急网络研究深度强化学习路由，提高路径对网络变化的适应能力\cite{drlrr}。这些工作说明，能量和时延可以直接参与路由决策，而不只是事后统计指标。

路径确定后，资源分配影响链路能够提供的速率和服务时间。Mozaffari 等从部署、信道和能量效率等方面讨论无人机通信设计，并总结优化与机器学习等求解途径\cite{mozaffari2019tutorial}。对于模型和环境信息较为明确的问题，可以依据速率与功率关系构建优化目标；在难以准确获得状态转移时，可以通过交互学习动作价值。前一类方法便于表达约束，后一类方法可利用运行数据适应变化。两者都需要明确有限电量如何分配到连续到达的通信需求，而不能仅根据当前链路速率决定发送强度。

当业务到达和未来信道未知时，队列方法提供了另一种在线分配依据。Tassiulas 和 Ephremides 用队列积压描述尚未满足的传输需求，研究多跳网络的吞吐量与稳定性\cite{tassiulas1992stability}；Neely 的李雅普诺夫优化方法进一步将长期资源约束转化为当前决策中的代价\cite{neely2010stochastic}。反压路由依据相邻节点的队列差分配转发量，Ying 等将其与最短路径结合以减少绕行\cite{related2}，Yu 和 Neely 分析了联合速率控制与路由中的效用和队列长度\cite{algo1}。在无人机系统中，Lin 等利用队列和能量状态安排太阳能无人机的资源分配与轨迹\cite{intro2}。这些研究表明，队列可用于联系当前负载与长期资源使用，但具体动作和约束仍由应用条件决定。

对本文的连续观测场景，路由与功率控制需要同时响应两个变化：一是当前数据积压和链路速率，二是关键节点的剩余电量及此前消耗。降低一次传输的能耗有助于节能，但持续使用同一中继仍可能缩短服务时间；均匀轮换中继也可能因使用较差链路而降低吞吐量。本文据此将本地队列、链路代价和累计能耗共同用于分布式决策，研究有限通信电量下的转发能力与网络可用时间。

\subsection{协同观测中的消息选择}

多节点协同观测利用不同视角补充单个节点看不到的信息，但节点增加也会带来更多待传数据。消息选择研究的基本问题是：在允许交换的数据量有限时，确定哪些节点和哪些内容最有助于接收端完成判断。无人机和车辆协同感知都存在局部视野不完整、观测重叠和无线带宽有限的问题，因而其信息筛选方法与本文直接相关；具体选择对象则随图像、特征或文本等消息形式而变化。

早期协同感知方法主要研究如何融合其他节点的信息。车车通信网络 V2VNet 交换压缩的中间特征，用于联合感知与预测\cite{v2vnet}；DiscoNet 通过知识蒸馏学习协作关系，对来自不同节点的特征分配权重\cite{disconet}。这些方法说明了局部观测之间的互补价值，但融合更多特征也需要更多传输资源。因此，后续工作将通信对象和发送内容纳入决策，而不再默认所有节点均提供完整特征。

Who2com 通过可学习的请求与匹配机制选择信息提供者\cite{who2com}，When2com 进一步根据观测需求决定通信时机和协作关系\cite{when2com}。Where2comm 利用空间置信图筛选对感知结果重要的稀疏区域，并随可用带宽调整交换范围，其实验同时涉及车辆和无人机协同感知\cite{where2comm}。How2comm 从空间位置和特征通道两个维度减少冗余，并补偿传输延迟造成的特征不同步\cite{how2comm}。这一发展过程体现出明确的变化：通信选择从确定信息来源，逐步细化到来源中具体需要发送的部分。

消息价值还与先前的通信有关。接收端已经获得某一区域的信息后，再发送相似内容的作用可能下降；一个单独使用时不够可靠的视角，也可能补充其他视角缺少的证据。学习式通信调度可以根据任务反馈决定是否继续交换信息\cite{ic3net,schednet}。本文采用中心选择消息的方式，借鉴的是这种根据任务状态决定发送时机的思路，并将接收端已有判断、来源质量和已选历史作为选择依据。

现有空间筛选方法主要评价图像区域或中间特征，而类别判断、陈述和摘要具有不同的长度与信息密度。对这类候选，仅按单条消息的置信度排序，容易忽略其与已有内容的重复程度。本文第二项研究因此同时选择消息来源和表达粒度，并在每次选择后更新候选评价，在消息总长度约束内提高中心判断的准确性。

\subsection{面向任务的语义通信}

消息选择需要回答信息对任务有多大作用。经典通信理论通过熵、信道容量和率失真描述有限资源下的信息传输与恢复\cite{shannon1948,cover2006information}。对于协同判断，接收端通常只需要观测中与当前问题相关的部分。信息瓶颈方法通过限制中间表示保留的输入信息、维持其与目标的相关性，为任务相关的信息压缩提供了分析依据\cite{tishby2000ib}。面向任务的语义通信据此将最终判断或推理结果纳入通信评价\cite{semantic-communications-jsac,zhou2024goal}。

深度联合信源信道编码通过联合训练编码器和解码器，学习适应信道的传输表示。Bourtsoulatze 等研究图像到信道符号的映射，改善受限带宽和噪声条件下的图像传输\cite{bourtsoulatze2019deepjscc}。深度语义通信模型（Deep Semantic Communication，DeepSC）利用注意力网络提取并传输文本含义\cite{xie2021deepsc,vaswani2017attention}；Shao 等面向边缘推理，以信息瓶颈保留有效特征，并采用可变长度表示适应信道\cite{shao2021task}。这些方法说明，保留任务所需信息能够减少对完整数据恢复的依赖。用于多源观测时，还需进一步确定应当编码和发送哪些来源，才能避免对重复内容进行高效但无益的传输。

视觉语言模型（Vision-Language Model，VLM）为构造显式观测消息提供了工具。对比语言图像预训练模型（Contrastive Language-Image Pre-training，CLIP）通过图文学习建立共享表示\cite{clip}；视觉指令模型能够根据图像和问题生成回答或描述\cite{blip2,llava}。局部观测由此可以表示为类别判断、置信度及简短文本，供接收中心融合使用。模型输出仍可能受遮挡和视角偏差影响，因此生成消息后还需评价其来源及其对当前判断的补充作用。

任务查询决定需要保留的信息。视觉问答（Visual Question Answering，VQA）研究根据问题定位相关区域\cite{where-to-look-vqa}；视觉问答数据集第二版（VQAv2）利用同一问题对应不同答案的图像对减少语言先验偏差\cite{vqa-v2}，面向真实场景的组合问答数据集 GQA 则构造多步推理问题\cite{gqa}。常见物体上下文数据集（Common Objects in Context，COCO）的图像与标注可用于形成局部观测\cite{coco}。本文利用图像与任务查询生成候选消息，并保持提取模型参数固定，使不同通信策略面对相同内容，比较内容选择本身对判断结果和通信量的影响。

\subsection{语义消息的可靠传输与重传控制}

消息经过选择后，仍可能因链路丢失而未被接收。ARQ 依据接收确认和超时安排补发，停等和选择重传分别在协议复杂度、并发能力与资源使用方面作出不同取舍\cite{lin1984arq}。混合自动重传请求（Hybrid Automatic Repeat Request，HARQ）进一步结合差错控制编码与多轮重传。对观测任务而言，重传的收益不仅取决于分组是否缺失，还取决于缺失内容是否会影响最终判断；因此，传输可靠性需要与任务结果联系起来。

语义重传研究已经开始利用任务质量决定补充哪些信息。Sheng 等针对协同感知提取重要特征，结合语义错误检测与 HARQ 改善衰落信道下的感知表现\cite{sheng2024harq}。Li 等研究重复语义信号、多轮接收融合以及增加语义冗余的重传方式\cite{li2025semharq}。这些方法通过任务相关内容和多轮融合提高重传效率，说明传输失败后并非只能按原有字节顺序重复发送全部数据。

重传内容确定后，还存在发送时机的问题。确认消息迟到或丢失时，发送端无法直接区分“接收端缺少该内容”和“内容已到达但尚未获知”两种状态。前一种情况可能需要补发，后一种情况则可能因补发消耗本可用于其他消息的资源。Huang 等用信息效用损失描述收发状态不一致对传输决策的影响，并在卫星融合网络中联合安排采样、编码和重传\cite{uoi2024satellite}。这说明任务收益评价需要使用接收端状态，而不仅是待发消息自身的重要性。

部分可观测马尔可夫决策过程（Partially Observable Markov Decision Process，POMDP）用条件概率表示无法直接观察的状态，并结合历史信息选择动作\cite{kaelbling1998pomdp}。对于观测消息，可以利用发送成功概率、发送记录和已返回的确认估计接收集合，再计算新发或重传的预期收益。考虑更多未来动作有助于分析等待的作用，但会增加在线计算量；模型预测控制通过限制预测时域、按新观测重复求解控制这一开销\cite{mayne2000mpc}。本文第三项研究据此比较有限步内的新发、重传、等待和停止，并将剩余任务时间与通信电量纳入决策。

\subsection{现有研究的不足}

上述工作分别改善了网络传输、信息筛选和重传效率。用于通信资源受限的多无人机协同观测时，仍需把这些能力与连续任务中的实际资源支出联系起来，主要存在以下三个问题。

（1）当前传输收益与后续网络可用时间之间需要进一步协调。能量感知和学习式路由已经将剩余电量、时延及投递结果纳入决策\cite{qtar,drlfr,drlrr}，但同一条高收益路径被持续使用后，关键中继的累计消耗可能改变其后续服务能力。只优化瞬时速率，容易集中消耗这些节点；只按剩余电量分流，又可能牺牲有效转发量。本文关注本地业务积压、链路速率与累计能耗共同作用时的下一跳和功率选择，利用长期能量反馈协调当前转发需求与后续通信能力。

（2）单条消息质量不足以确定其在当前已收信息之后的价值。现有语义表示和协同感知方法能够减少冗余数据\cite{xie2021deepsc,where2comm,how2comm}，但不同来源、不同表达长度的显式消息还具有相互补充或重复的关系。高置信度可能来自受遮挡的视角，较长表达也未必补充中心缺少的判断依据；每次选择还会改变剩余预算与后续消息的作用。本文因此在统一的消息长度预算内，根据来源可靠性、中心判断和通信历史更新候选评价，同时学习消息选择与停止时机。

（3）内容重要性与当前重传必要性需要联合判断。语义重传能够优先补充任务相关信息\cite{sheng2024harq,li2025semharq}，但在确认不及时的条件下，一条重要消息也可能已经到达。发送端需要先估计接收状态，再判断立即补发或等待确认是否更有利。此外，重复消息即使提高了模型置信度，也可能不改变最终类别，因而不应仅按置信度增长估计任务收益。本文第三项研究结合接收状态估计、同类别消息组合的收益一致性和等待决策，减少低收益发送，并把当前节省的能量保留给后续任务。

这三个不足对应观测信息回传中的连续环节。路由研究改善可用传输条件，消息选择研究确定有限资源应承载的内容，发送与重传研究根据实际接收情况调整后续通信。本文以这一关系组织方法设计，并通过最终接收信息是否支持按时正确判断，说明各项研究与任务可靠完成之间的联系。


\section{本文研究内容及章节安排}
\subsection{主要研究内容}

围绕保障持续回传、筛选有效语义与实现按时可靠交付三个核心环节，本文构建了如图~\ref{fig:intro:framework} 所示的整体研究框架。三项研究立足于同一协同观测任务，通过消息负载、链路速率与通信能耗紧密衔接，主要内容如下。

（1）\ResearchOneTitle{}。针对高动态无人机网络中机载通信电量有限、转发效率与网络持续服务时间难以兼顾的问题，本文提出了基于李雅普诺夫优化的能量感知分布式路由与功率控制方法（Lyapunov-based Energy-Aware Decentralized Routing，LEADR）。该方法将数据队列积压、节点累计能耗与虚拟能量队列共同纳入逐时隙优化目标，使各节点能够结合当前信道与干扰状态直接求解最优发射功率，并综合邻居剩余电量与到中心的地理进度挑选下一跳，同时依据节点间的相对运动预测自适应调整邻居探测周期。LEADR 仅依赖本地状态即可协调当前转发负载与长期耗电，在维持较高转发吞吐量与低丢包率的同时避免关键中继提前耗尽电量。

（2）\ResearchTwoTitle{}。针对链路连通后单次任务消息长度预算受限、多机视角重叠与遮挡导致冗余或错误描述挤占传输资源的问题，本文提出了基于强化学习的语义消息选择方法 EviComRL。在参数固定的视觉语言模型离线生成多视角、不同详略程度候选消息的基础上，该方法先通过结合开销的初筛缩小候选规模，再由策略网络综合发送方观测可靠性、中心当前判断、已选历史与链路代价动态评估剩余候选，通过屏蔽已选或超预算动作并采用近端策略优化算法训练，自主学习选择哪些消息以及何时停止传输，从而在较低词元开销下提升地面中心的判断准确率。

（3）\ResearchThreeTitle{}。针对候选消息在无线回传中同时面临信道丢包与确认反馈延迟或丢失、盲目重传耗电而过度等待易致超时的问题，本文提出了一致性感知的信念反馈控制方法（Consistency-aware Belief Feedback Control，C-BFC）。该方法首先训练轻量级神经网络估计不同接收消息组合的判断准确程度，并将预测类别相同的组合估值取均值以保持一致，避免诱发不改变最终类别判断的冗余发送；在线传输时，控制器依据已到达的累计确认和后续发送记录推算接收端当前收到各条消息的概率，利用跨任务虚拟能量队列控制能耗支出，并在剩余时限内权衡两步发送前瞻与等待确认的预期回报，动态决定发送新消息、重传、等待或停止，以较低能耗提升任务按时正确完成率。

\begin{figure}[!htb]
\centering
% Read study 1 first, then studies 2 and 3 from left to right.
% Dashed arrows show functional relations; solid inter-study arrows show
% candidate messages or calculations already reused in the experiments.
\begingroup
\begin{tikzpicture}[
  x=1mm,y=1mm,>=Stealth,
  font=\fontsize{8.7bp}{10.9bp}\selectfont,
  rfborder/.style={draw=black,line width=.5pt},
  rfbox/.style={rfborder,fill=white,align=center,inner sep=.6mm,outer sep=0pt},
  rfhead/.style={align=center,text width=61mm,inner sep=0pt,
    font=\fontsize{9bp}{11.4bp}\selectfont\bfseries},
  rfbody/.style={align=left,anchor=north west,text width=61mm,
    inner sep=0pt,outer sep=0pt},
  rfmethod/.style={align=center,inner sep=0pt,
    font=\fontsize{10bp}{12bp}\selectfont\bfseries},
  rflinktext/.style={align=left,inner sep=0pt,
    font=\fontsize{8.4bp}{10.6bp}\selectfont},
  rfedge/.style={-{Stealth[length=1.5mm,width=1.05mm]},
    draw=black,line width=.6pt},
  rffunction/.style={rfedge,line width=.55pt,dashed,
    dash pattern=on 1.4mm off .8mm}
]
\path[use as bounding box] (0,0) rectangle (148,-145);

\node[rfbox,text width=145mm,minimum height=7mm,
  font=\fontsize{10.2bp}{12.8bp}\selectfont\bfseries]
  at (74,-3.5) {\ThesisTitle};

\draw[rfborder] (0,-10) rectangle (148,-27);
\node[align=left,text width=144mm,inner sep=0pt,
  font=\fontsize{8.9bp}{11.4bp}\selectfont] at (74,-18.5)
  {\textbf{应用场景：}多无人机回传不同视角的局部信息，由地面中心汇集信息并判断现场状况。\\
   \textbf{共同约束：}通信电量有限，每项任务允许发送的消息总长度和完成时间受限。\\
   \textbf{总体目标：}在时限内汇集有效信息并正确判断，以较低通信能耗完成连续观测任务。};
\draw[rfedge] (74,-27) -- (74,-32);

% Study 1: routing and power are distinct but coupled control variables.
\draw[rfborder] (0,-32) rectangle (148,-64);
\node[rfhead,text width=144mm] at (74,-36)
  {研究问题1：\ResearchOneTitle};
\draw[rfborder,line width=.35pt] (2,-39.5) -- (146,-39.5);
\node[align=center,text width=144mm,inner sep=0pt,
  font=\fontsize{8.7bp}{10.9bp}\selectfont\bfseries] at (74,-43)
  {目标：在通信能量约束下最大化网络聚合吞吐量，并延长网络可用时间。};
\draw[rfborder,line width=.35pt] (48,-46.5) -- (48,-62);
\draw[rfborder,line width=.35pt] (98,-46.5) -- (98,-62);
\node[rfbody,text width=43mm] at (3,-47)
  {\textbf{具体问题}\\
   功率偏低导致积压与丢包，\\
   高功率及集中转发加快耗电，\\
   难以兼顾吞吐量与使用时间。};
\node[rfbody,text width=45mm] at (51,-47)
  {\textbf{方法与输出（LEADR）}\\
   联合确定下一跳和发射功率；\\
   按积压与能耗调整转发强度；\\
   随节点运动更新邻居信息。};
\node[rfbody,text width=44mm] at (101,-47)
  {\textbf{创新点}\\
   在吞吐量优化中引入累计能耗，\\
   抑制持续高耗能发送；\\
   按邻居电量修正转发权重。};

% The two directions describe a functional dependency, not online integration.
\draw[rffunction] (10,-64) -- (10,-80);
\node[rflinktext,anchor=west,text width=47mm] at (14,-69.5)
  {传输速率与发送耗能};
\draw[rffunction] (62,-80) -- (62,-64);
\node[rflinktext,anchor=east,text width=47mm,align=right] at (58,-76)
  {消息数量与长度形成转发负载};
% Study 3 reuses the link-rate and energy calculations, not the LEADR policy.
\draw[rfedge] (99,-64) -- (99,-80);
\node[rflinktext,anchor=west,text width=44mm] at (102,-72)
  {LEADR 的速率与能耗模型\\计算发送时间与能耗};

% Studies 2 and 3 follow in left-to-right reading order.
\draw[rfborder] (0,-80) rectangle (65,-140);
\draw[rfborder] (83,-80) rectangle (148,-140);
\node[rfhead] at (32.5,-86.6)
  {研究问题2：\\\ResearchTwoTitle};
\node[rfhead] at (115.5,-86.6)
  {研究问题3：\\\ResearchThreeTitle};
\node[rfmethod] at (32.5,-96) {EviComRL};
\node[rfmethod] at (115.5,-96) {C-BFC};
\draw[rfborder,line width=.35pt] (2,-98.9) -- (63,-98.9);
\draw[rfborder,line width=.35pt] (85,-98.9) -- (146,-98.9);

\node[rfbody] at (2,-101)
  {\textbf{具体问题：}消息总长度受限；\\
   重复或不可靠的观测占用发送机会，\\
   中心缺少的有效信息可能未被选中。};
\node[rfbody] at (2,-113.1)
  {\textbf{目标：在有限通信量下提高判断准确率。}};
\node[rfbody] at (2,-118)
  {\textbf{方法：}按观测质量与发送代价筛选消息，\\学习消息来源、长度和停止时机。};
\node[rfbody] at (2,-126.3)
  {\textbf{创新点：}结合来源可靠性与选择历史，\\随中心判断更新候选评价，\\减少重复或低质量内容的发送。};

\node[rfbody] at (85,-101)
  {\textbf{具体问题：}确认延迟或丢失时，\\
   发送端难以确定已收消息；\\
   重复补发耗电，持续等待可能超时。};
\node[rfbody] at (85,-113.1)
  {\textbf{目标：提高按时正确完成率并控制能耗。}};
\node[rfbody] at (85,-118)
  {\textbf{方法：}按发送记录与确认估计接收状态；\\
   比较继续发送、等待和停止的收益。};
% This equality applies to correctness of nonempty sets within one task.
\node[rfbody] at (85,-126.3)
  {\textbf{创新点：}同一任务中，相同类别的\\非空消息组合采用一致收益；\\
   结合接收概率评价重传与等待。};

\draw[rfedge] (65,-114.5) -- (83,-114.5);
\node[rflinktext,align=center,text width=17mm] at (74,-108.5)
  {候选消息\\及其长度};
\node[align=center,inner sep=0pt,text width=148mm,
  font=\fontsize{8.2bp}{10.3bp}\selectfont] at (74,-143.3)
  {虚线表示传输条件与负载的联系；实线表示候选消息与链路计算模型的使用关系};
\end{tikzpicture}
\endgroup

\bicaption{本文整体研究框架}{Overall research framework of this thesis}
\label{fig:intro:framework}
\end{figure}

\subsection{整体研究方案}

首先，为了在统一框架下分析通信能量、消息长度和任务时限如何影响最终的协同观测结果，本文建立了多无人机协同观测的系统通信模型，明确动态无线链路、数据队列与电量演化、三类资源约束以及任务按时正确完成条件，为后续各部分算法设计提供统一依据，具体描述见第~\ref{ch:foundations} 章。

在此基础上，围绕多跳网络持续回传的问题，研究一设计了分布式路由与功率控制方法，负责确定各无人机节点的转发路径与发射强度，具体描述见第~\ref{rq1:chapter} 章。各节点依据本地数据队列积压、实时信道状态与邻居剩余电量，在线输出下一跳中继节点与对应的发射功率。这一输出直接决定了无线链路的传输速率和单次发送耗电，为上层评估消息传输的时间与能量开销提供了底层参数。

接着，针对链路连通后单次任务消息长度有限、多机视角重叠且存在遮挡的问题，研究二设计了语义消息选择方法，具体描述见第~\ref{rq2:chapter} 章。该方法利用底层链路速率与传输开销，结合发送节点的观测可靠程度和地面中心当前判断，从各视角不同详略程度的候选中筛选有效消息并决定停止时机，最终输出包含来源节点、语义内容和消息长度的待发候选集合。此外，筛选出的消息数量与文本长度构成了底层网络实际需要转发的数据负载，会直接影响各节点的队列积压与中继耗电速度。图~\ref{fig:intro:framework} 中的虚线即表示研究一与研究二在链路传输条件和业务负载上的相互影响。

最后，为了使筛选出的候选消息在无线丢包与确认反馈延迟或丢失的情况下按时送达地面中心，研究三设计了发送与重传控制方法，具体描述见第~\ref{rq3:chapter} 章。该方法直接接收研究二输出的候选消息作为输入，同时沿用研究一的链路速率与发送能耗计算方式，由消息长度与链路速率算出单次发送时长，再结合发射功率核算相应耗电，通过一致的数据流与通信参数将前两项研究连接起来，对应图~\ref{fig:intro:framework} 中的实线流程。实际传输时，控制器结合历史发送记录、地面中心返回的累计确认、节点剩余电量与任务剩余时间，在线决定发送新消息、重传、等待还是停止，最终确定地面中心在截止时间前实际收到的消息集合，从而提高任务按时正确完成率。

\FloatBarrier
\subsection{章节安排}

本文其余章节按如下结构组织。第~\ref{ch:foundations} 章建立多无人机协同观测的系统通信模型，明确定义通信能量、消息长度与任务时限约束以及任务按时正确完成条件，并介绍后续章节用到的优化与决策方法基础。第~\ref{rq1:chapter} 章至第~\ref{rq3:chapter} 章依次展开 LEADR 路由与功率控制、EviComRL 语义消息选择以及 C-BFC 发送与重传控制的研究，每章均包含问题建模、算法设计、实验设置与结果分析。最后，结论部分对全文的主要研究工作和实验结果进行总结。